% ====================================================================
% EXEMPLO.TEX — COMO CRIAR SEÇÕES, SUBSEÇÕES E SUBSUBSEÇÕES
% ====================================================================
% Este arquivo é só um EXEMPLO de uso, não é chamado pelo main.tex.
% Para usar um destes arquivos de verdade no corpo do documento, dê
% % ====================================================================
% EXEMPLO.TEX — COMO CRIAR SEÇÕES, SUBSEÇÕES E SUBSUBSEÇÕES
% ====================================================================
% Este arquivo é só um EXEMPLO de uso, não é chamado pelo main.tex.
% Para usar um destes arquivos de verdade no corpo do documento, dê
% % ====================================================================
% EXEMPLO.TEX — COMO CRIAR SEÇÕES, SUBSEÇÕES E SUBSUBSEÇÕES
% ====================================================================
% Este arquivo é só um EXEMPLO de uso, não é chamado pelo main.tex.
% Para usar um destes arquivos de verdade no corpo do documento, dê
% % ====================================================================
% EXEMPLO.TEX — COMO CRIAR SEÇÕES, SUBSEÇÕES E SUBSUBSEÇÕES
% ====================================================================
% Este arquivo é só um EXEMPLO de uso, não é chamado pelo main.tex.
% Para usar um destes arquivos de verdade no corpo do documento, dê
% \input{secoes/exemplo.tex} dentro de main.tex (onde hoje está
% \input{secoes/teste.tex}).
%
% Estrutura hierárquica disponível:
%   \section{...}       -> nível 1 (numerado 1, 2, 3...)
%   \subsection{...}    -> nível 2 (numerado 1.1, 1.2...)
%   \subsubsection{...} -> nível 3 (numerado 1.1.1, 1.1.2...)
% ====================================================================

\section{Nome da Seção}

Texto de abertura da seção. Todo o conteúdo escrito aqui, antes da
primeira \verb|\subsection|, é considerado parte da introdução desta
seção.

\subsection{Nome da Subseção}

Texto da subseção. Uma seção pode ter quantas subseções forem
necessárias, bastando repetir o comando \verb|\subsection{...}|.

\subsubsection{Nome da Subsubseção}

Texto da subsubseção — o menor nível de divisão usado neste modelo.
Uma subseção pode ter quantas subsubseções forem necessárias.

\subsubsection{Outra Subsubseção}

Basta repetir \verb|\subsubsection{...}| para criar quantas forem
necessárias dentro da mesma subseção.

\subsection{Outra Subseção}

E o mesmo vale para subseções: repita \verb|\subsection{...}| dentro
da seção quantas vezes precisar.

% --------------------------------------------------------------------
% COMO INSERIR UMA IMAGEM
% --------------------------------------------------------------------
% O comando \includegraphics (do pacote graphicx, já carregado em
% includes.tex) insere uma imagem no documento. Aqui usamos o ambiente
% "figure" para que a imagem fique centralizada e com uma legenda
% numerada automaticamente (Figura 1, Figura 2, ...).
%
%   [h]                  -> pede ao LaTeX para tentar posicionar a
%                           imagem "aqui mesmo" (here), no ponto do
%                           texto onde o comando foi escrito.
%   \centering           -> centraliza a imagem horizontalmente.
%   width=10cm           -> define a largura da imagem em centímetros
%                           (ajuste esse valor conforme o tamanho
%                           desejado; a altura é ajustada automaticamente
%                           para manter a proporção original da imagem).
%   \caption{...}        -> texto da legenda, exibido abaixo da imagem.
%   \label{fig:exemplo}  -> permite referenciar a imagem em outro ponto
%                           do texto com \ref{fig:exemplo} (ex.: "ver
%                           Figura~\ref{fig:exemplo}").
%
% O caminho da imagem é relativo à raiz do projeto (onde está o
% main.tex), então "imag/imag-exemplo.png" busca o arquivo dentro da
% pasta "imag/".
\begin{figure}[H]
    \centering
    \includegraphics[width=4cm]{imag/imag-exemplo.jpg}
    \caption{Exemplo de imagem inserida no documento.}
    \label{fig:exemplo}
\end{figure}

\section{Outra Seção}

Para começar uma nova seção (nível 1), basta usar \verb|\section{...}|
novamente. Não é necessário fechar nada antes. dentro de main.tex (onde hoje está
% \input{secoes/teste.tex}).
%
% Estrutura hierárquica disponível:
%   \section{...}       -> nível 1 (numerado 1, 2, 3...)
%   \subsection{...}    -> nível 2 (numerado 1.1, 1.2...)
%   \subsubsection{...} -> nível 3 (numerado 1.1.1, 1.1.2...)
% ====================================================================

\section{Nome da Seção}

Texto de abertura da seção. Todo o conteúdo escrito aqui, antes da
primeira \verb|\subsection|, é considerado parte da introdução desta
seção.

\subsection{Nome da Subseção}

Texto da subseção. Uma seção pode ter quantas subseções forem
necessárias, bastando repetir o comando \verb|\subsection{...}|.

\subsubsection{Nome da Subsubseção}

Texto da subsubseção — o menor nível de divisão usado neste modelo.
Uma subseção pode ter quantas subsubseções forem necessárias.

\subsubsection{Outra Subsubseção}

Basta repetir \verb|\subsubsection{...}| para criar quantas forem
necessárias dentro da mesma subseção.

\subsection{Outra Subseção}

E o mesmo vale para subseções: repita \verb|\subsection{...}| dentro
da seção quantas vezes precisar.

% --------------------------------------------------------------------
% COMO INSERIR UMA IMAGEM
% --------------------------------------------------------------------
% O comando \includegraphics (do pacote graphicx, já carregado em
% includes.tex) insere uma imagem no documento. Aqui usamos o ambiente
% "figure" para que a imagem fique centralizada e com uma legenda
% numerada automaticamente (Figura 1, Figura 2, ...).
%
%   [h]                  -> pede ao LaTeX para tentar posicionar a
%                           imagem "aqui mesmo" (here), no ponto do
%                           texto onde o comando foi escrito.
%   \centering           -> centraliza a imagem horizontalmente.
%   width=10cm           -> define a largura da imagem em centímetros
%                           (ajuste esse valor conforme o tamanho
%                           desejado; a altura é ajustada automaticamente
%                           para manter a proporção original da imagem).
%   \caption{...}        -> texto da legenda, exibido abaixo da imagem.
%   \label{fig:exemplo}  -> permite referenciar a imagem em outro ponto
%                           do texto com \ref{fig:exemplo} (ex.: "ver
%                           Figura~\ref{fig:exemplo}").
%
% O caminho da imagem é relativo à raiz do projeto (onde está o
% main.tex), então "imag/imag-exemplo.png" busca o arquivo dentro da
% pasta "imag/".
\begin{figure}[H]
    \centering
    \includegraphics[width=4cm]{imag/imag-exemplo.jpg}
    \caption{Exemplo de imagem inserida no documento.}
    \label{fig:exemplo}
\end{figure}

\section{Outra Seção}

Para começar uma nova seção (nível 1), basta usar \verb|\section{...}|
novamente. Não é necessário fechar nada antes. dentro de main.tex (onde hoje está
% \input{secoes/teste.tex}).
%
% Estrutura hierárquica disponível:
%   \section{...}       -> nível 1 (numerado 1, 2, 3...)
%   \subsection{...}    -> nível 2 (numerado 1.1, 1.2...)
%   \subsubsection{...} -> nível 3 (numerado 1.1.1, 1.1.2...)
% ====================================================================

\section{Nome da Seção}

Texto de abertura da seção. Todo o conteúdo escrito aqui, antes da
primeira \verb|\subsection|, é considerado parte da introdução desta
seção.

\subsection{Nome da Subseção}

Texto da subseção. Uma seção pode ter quantas subseções forem
necessárias, bastando repetir o comando \verb|\subsection{...}|.

\subsubsection{Nome da Subsubseção}

Texto da subsubseção — o menor nível de divisão usado neste modelo.
Uma subseção pode ter quantas subsubseções forem necessárias.

\subsubsection{Outra Subsubseção}

Basta repetir \verb|\subsubsection{...}| para criar quantas forem
necessárias dentro da mesma subseção.

\subsection{Outra Subseção}

E o mesmo vale para subseções: repita \verb|\subsection{...}| dentro
da seção quantas vezes precisar.

% --------------------------------------------------------------------
% COMO INSERIR UMA IMAGEM
% --------------------------------------------------------------------
% O comando \includegraphics (do pacote graphicx, já carregado em
% includes.tex) insere uma imagem no documento. Aqui usamos o ambiente
% "figure" para que a imagem fique centralizada e com uma legenda
% numerada automaticamente (Figura 1, Figura 2, ...).
%
%   [h]                  -> pede ao LaTeX para tentar posicionar a
%                           imagem "aqui mesmo" (here), no ponto do
%                           texto onde o comando foi escrito.
%   \centering           -> centraliza a imagem horizontalmente.
%   width=10cm           -> define a largura da imagem em centímetros
%                           (ajuste esse valor conforme o tamanho
%                           desejado; a altura é ajustada automaticamente
%                           para manter a proporção original da imagem).
%   \caption{...}        -> texto da legenda, exibido abaixo da imagem.
%   \label{fig:exemplo}  -> permite referenciar a imagem em outro ponto
%                           do texto com \ref{fig:exemplo} (ex.: "ver
%                           Figura~\ref{fig:exemplo}").
%
% O caminho da imagem é relativo à raiz do projeto (onde está o
% main.tex), então "imag/imag-exemplo.png" busca o arquivo dentro da
% pasta "imag/".
\begin{figure}[H]
    \centering
    \includegraphics[width=4cm]{imag/imag-exemplo.jpg}
    \caption{Exemplo de imagem inserida no documento.}
    \label{fig:exemplo}
\end{figure}

\section{Outra Seção}

Para começar uma nova seção (nível 1), basta usar \verb|\section{...}|
novamente. Não é necessário fechar nada antes. dentro de main.tex (onde hoje está
% \input{secoes/teste.tex}).
%
% Estrutura hierárquica disponível:
%   \section{...}       -> nível 1 (numerado 1, 2, 3...)
%   \subsection{...}    -> nível 2 (numerado 1.1, 1.2...)
%   \subsubsection{...} -> nível 3 (numerado 1.1.1, 1.1.2...)
% ====================================================================

\section{Nome da Seção}

Texto de abertura da seção. Todo o conteúdo escrito aqui, antes da
primeira \verb|\subsection|, é considerado parte da introdução desta
seção.

\subsection{Nome da Subseção}

Texto da subseção. Uma seção pode ter quantas subseções forem
necessárias, bastando repetir o comando \verb|\subsection{...}|.

\subsubsection{Nome da Subsubseção}

Texto da subsubseção — o menor nível de divisão usado neste modelo.
Uma subseção pode ter quantas subsubseções forem necessárias.

\subsubsection{Outra Subsubseção}

Basta repetir \verb|\subsubsection{...}| para criar quantas forem
necessárias dentro da mesma subseção.

\subsection{Outra Subseção}

E o mesmo vale para subseções: repita \verb|\subsection{...}| dentro
da seção quantas vezes precisar.

% --------------------------------------------------------------------
% COMO INSERIR UMA IMAGEM
% --------------------------------------------------------------------
% O comando \includegraphics (do pacote graphicx, já carregado em
% includes.tex) insere uma imagem no documento. Aqui usamos o ambiente
% "figure" para que a imagem fique centralizada e com uma legenda
% numerada automaticamente (Figura 1, Figura 2, ...).
%
%   [h]                  -> pede ao LaTeX para tentar posicionar a
%                           imagem "aqui mesmo" (here), no ponto do
%                           texto onde o comando foi escrito.
%   \centering           -> centraliza a imagem horizontalmente.
%   width=10cm           -> define a largura da imagem em centímetros
%                           (ajuste esse valor conforme o tamanho
%                           desejado; a altura é ajustada automaticamente
%                           para manter a proporção original da imagem).
%   \caption{...}        -> texto da legenda, exibido abaixo da imagem.
%   \label{fig:exemplo}  -> permite referenciar a imagem em outro ponto
%                           do texto com \ref{fig:exemplo} (ex.: "ver
%                           Figura~\ref{fig:exemplo}").
%
% O caminho da imagem é relativo à raiz do projeto (onde está o
% main.tex), então "imag/imag-exemplo.png" busca o arquivo dentro da
% pasta "imag/".
\begin{figure}[H]
    \centering
    \includegraphics[width=4cm]{imag/imag-exemplo.jpg}
    \caption{Exemplo de imagem inserida no documento.}
    \label{fig:exemplo}
\end{figure}

\section{Outra Seção}

Para começar uma nova seção (nível 1), basta usar \verb|\section{...}|
novamente. Não é necessário fechar nada antes.